\scanpage{269}
\mspage{81}
% archive sheet number: 269
Donc

\[
f^*(G)\simeq(j\bar f)^*(\Phi)|I
\]

Or la flèche

\[
\alpha:\mathrm{id}_K\longrightarrow j\bar f,\quad
\deleted{\unsure{i\longrightarrow j\bar f(i)}}
\]

donne

\[
\underset{\substack{\parallel\\\Phi}}{\mathrm{id}_K^*(\Phi)}
\xleftarrow{\Phi(\alpha)}(j\bar f)^*(\Phi)
\qquad\text{fonctoriel en }\Phi
\]

d'où, en prenant la restriction à $I$

\[
i^*(\Phi)\overset{\mathrm{def}}{=}F\xleftarrow{u}f^*G.
\]

En fait, sans tellement se tortiller, on devrait pouvoir montrer que \deleted{\unsure{en}} pour \unsure{toute} catégorie $\mathcal M$, la catégorie $\underline{\operatorname{Hom}}(K^{\unsure{\mathrm o|\circ}},\mathcal M)$ des $\mathcal M$-préfaisceaux sur $K$ est \deleted{équivalente}\inserted{isomorphe} à celle des triples $(F,G,u)$, avec $F\in\mathrm{Ob}\,\underline{\operatorname{Hom}}(I^{\mathrm o},\mathcal M)$, $G\in\underline{\operatorname{Hom}}(J^{\mathrm o},\mathcal M)$, et $u:f^*(G)\longrightarrow F$. \deleted{\unsure{Mais}} \unsure{Ainsi}, $K$ apparaît comme solution d'un pb \deleted{\illegible{1}} universel

\deleted{\unsure{rel. au foncteur covariant}}

\[
\underline{\operatorname{Hom}}(K^{\mathrm o},\mathcal M)\longrightarrow
\{F,G,u\ \text{comme dessus}\}(\mathcal M)
\]

Mais pour y préciser la flèche $\longrightarrow$ de façon \unsure{conceptuelle|conceptuellement} simple, il est bon de se donner $K$ avec les structures supplémentaires

\diagram{\begin{tikzcd}[ampersand replacement=\&,row sep=tiny,column sep=small]
J\arrow[r,hook,"j"]\&K\arrow[l,bend left=35,"\bar f"]\&I\arrow[l,hook,no head,"i"']\\
\&\bar f j=\mathrm{id}_J\&
\end{tikzcd}
\quad$,\quad\mathrm{id}_K\xrightarrow{\alpha}j\bar f$.}

Et ce système
